% Regression: t2_finegrain — renders (a): if the transfer operator (quantum channel) E of the MPS.
% Formula: (a):  if the transfer operator (quantum channel) E of the MPS
% t2_finegrain — RMP arXiv:2011.12127, Fig. (fig:renorm2D), file
% fig2_fig2-pepe_newfig13.pdf: the inverse renormalization (fine-
% graining) step, panels (a) MPS and (b) PEPS.
%
% Formula (a):  if the transfer operator (quantum channel) E of the MPS
% tensor A is divisible, E = R.R, the site splits:
%     A  -->  A' A''    (one 3-leg tensor rewritten as two 3-leg
%                        tensors of the same bond dimension, one new
%                        virtual bond contracted, both physical legs
%                        kept).
% Formula (b):  the same step for a 2D PEPS tensor would need
%     A  -->  (2x2 block of tensors with bond dimension sqrt(D)),
% one 5-leg site tensor becoming a 2x2 sheared lattice patch.
%
% Declarative reading:
%   (a) two one-row grids (physical=up, default dot skin — the
%       reference draws dots) joined by the plain map arrow -->;
%   (b) two tenkzlattice pictures, frame=oblique (the sheared PEPS
%       sheet), boundary legs (the cut virtual bonds dangle), and
%       physical=up for the on-site leg; sizes 1x1 --> 2x2.
% The arrows are algebra between pictures, not tensor ink.
\documentclass[varwidth=19cm, border=6pt]{standalone}
\usepackage{amsmath, amssymb}
\usepackage{tenkz}
\begin{document}
\[
  (a)\quad
  % one MPS site: 3-leg dot
  \begin{tenkz}[physical=up]
    \tn{}
  \end{tenkz}
  \;\longrightarrow\;
  % split into two sites of the same bond dimension
  \begin{tenkz}[physical=up]
    \tn{} & \tn{}
  \end{tenkz}
  \qquad\qquad
  (b)\quad
  % one PEPS site: 5-leg tensor on the sheared sheet
  \begin{tenkzlattice}[rows=1, cols=1, boundary legs, physical=up,
                       frame=oblique]
  \end{tenkzlattice}
  \;\longrightarrow\;
  % the would-be sqrt(D) fine-graining: a 2x2 patch
  \begin{tenkzlattice}[rows=2, cols=2, boundary legs, physical=up,
                       frame=oblique]
  \end{tenkzlattice}
\]
\end{document}
